\frontmatterpage
\unchapter{Anotācija}

Šis bakalaura darbs pētīja informācijas sistēmu integrācijas briedumu projektu
vadības biroju (PVB) vidē Latvijas vidējos un lielos uzņēmumos. Tā kā organizācijas
arvien vairāk paļaujas uz vairākiem digitālajiem rīkiem projektu plānošanai,
pārskatošanai un pārvaldībai, sistēmu integrācijas pakāpei ir tieša ietekme uz
datu kvalitāti, pārskatu efektivitāti un pārvaldības standartizāciju. Neskatoties
uz tās praktisko nozīmi, šī tēma Latvijas kontekstā bija saņēmusi ierobežotu
empīrisku uzmanību.

Bakalaura darbs sastāv no piecām nodaļām. 1.\ nodaļā aplūkoti teorētiskie pamati,
definējot galvenos jēdzienus, piemēram, PVB, informācijas sistēmu integrāciju,
brieduma modeļus, biznesa inteliģenci un datu kvalitāti, kā arī apskatīti esošie
pētījumi par PVB digitalizāciju un integrācijas brieduma ietvariem. 2.\ nodaļā
sniegta pētījuma metodoloģija, iekļaujot četru līmeņu Integrācijas Brieduma
Modeļa izstrādi un strukturēta aptaujas instrumenta veidošanu. 3.\ nodaļā veikts
empīriskais pētījums, prezentējot aptaujas rezultātus, kas iegūti no Latvijas
uzņēmumu pārstāvjiem ar formālām PVB struktūrām, papildināts ar vienu strukturētu
ekspertu interviju un statistisko analīzi. 4.\ nodaļā rezultāti tika interpretēti
esošās literatūras kontekstā un izstrādāti pieci praktiski ieteikumi informācijas
sistēmu integrācijas brieduma uzlabošanai PVB vidē. 5.\ nodaļā aplūkoti nākotnes
pētījumu virzieni un plašākas sociālās un ekonomiskās sekas.

Pētījuma metodoloģiskais pamats ietver sistemātisku literatūras apskatu,
aptaujas bāzētu kvantitatīvo pētījumu, vienu strukturētu ekspertu interviju,
aprakstošo statistisko analīzi, kā arī Spīrmena ranga korelācijas analīzi.

Pētījums atklāja, ka 70\% PVB vižu Latvijā darbojas vidējā integrācijas brieduma
līmenī (2.\ vai 3.\ līmenis), un ka augstāks integrācijas briedums ir statistiski
nozīmīgi saistīts ar lielāku pārskatu automatizāciju (r=0,405, p=0,010), labāku
datu kvalitāti (r=0,321, p=0,043) un stiprāku pārvaldības standartizāciju
(r=0,403, p=0,010). Visi četri pētījuma hipotēzi tika apstiprināti. Rezultāti
sniedz uz pierādījumiem balstītus norādījumus organizācijām, kas vēlas uzlabot
savas projektu vadības informācijas sistēmas.

\thesisscopeLV

\textbf{Atslēgvārdi:} Projektu vadības birojs, Informācijas sistēmu integrācija,
Integrācijas brieduma modelis, Pārskatu automatizācija, Datu kvalitāte,
Biznesa inteliģence, Pārvaldības standartizācija.
